% ============================================================================
% LANGENTWURF LE-02d: Wiederholung + Stundenleistung Mi familia
% Spanisch Klasse 7 - Sequenz 2: Mi familia
% ============================================================================
% Tier: 1 (vollständig)
% Doppelstunde: 2d (Std. 7-8 von 8) - ABSCHLUSS
% Fachfarbe: #c41e3a (Karminrot)
% ============================================================================

\documentclass[11pt,a4paper]{article}

\usepackage[utf8]{inputenc}
\usepackage[T1]{fontenc}
\usepackage[ngerman]{babel}
\usepackage{geometry}
\usepackage{fancyhdr}
\usepackage{lastpage}
\usepackage{xcolor}
\usepackage{tcolorbox}
\usepackage{enumitem}
\usepackage{tabularx}
\usepackage{longtable}
\usepackage{array}
\usepackage{booktabs}
\usepackage{amssymb}

% ============================================================================
% KONFIGURATION
% ============================================================================

\newcommand{\fach}{Spanisch}
\newcommand{\fachfarbe}{c41e3a}
\newcommand{\jahrgangsstufe}{7}
\newcommand{\schulart}{Gesamtschule}
\newcommand{\stundenthema}{Wiederholung + Stundenleistung Mi familia}
\newcommand{\reihenthema}{Mi familia}
\newcommand{\stundennummer}{2d}
\newcommand{\reihenstunden}{4}
\newcommand{\gession}{A1}
\newcommand{\lehrwerk}{Qué pasa 1}
\newcommand{\lehrwerkseite}{S. 54-69}
\newcommand{\lerngruppengroesse}{24}

% ============================================================================
% LAYOUT
% ============================================================================
\geometry{left=2.5cm, right=2.5cm, top=2.5cm, bottom=2.5cm}
\definecolor{fach}{HTML}{\fachfarbe}
\definecolor{gray}{HTML}{666666}
\definecolor{lightgray}{HTML}{f5f5f5}
\definecolor{successgreen}{HTML}{2a7a4b}
\definecolor{warningorange}{HTML}{b35c00}
\definecolor{basis}{HTML}{2a7a4b}
\definecolor{standard}{HTML}{2c5aa0}
\definecolor{erweiterung}{HTML}{7b1fa2}

\tcbuselibrary{skins,breakable}

\newtcolorbox{infobox}[1][]{
    colback=lightgray,
    colframe=fach,
    fonttitle=\bfseries,
    title=#1,
    breakable
}

\newtcolorbox{scorebox}{
    colback=white,
    colframe=fach,
    fonttitle=\bfseries,
    title=Didaktik-Score,
    breakable
}

\pagestyle{fancy}
\fancyhf{}
\fancyhead[L]{\small\textcolor{gray}{\fach{} -- Jg. \jahrgangsstufe{} -- \stundenthema}}
\fancyhead[R]{\small\textcolor{gray}{LE-\stundennummer{} | Tier 1}}
\fancyfoot[C]{\small\thepage{} / \pageref{LastPage}}
\renewcommand{\headrulewidth}{0.4pt}

% Differenzierungsmarker
\newcommand{\B}{\textcolor{basis}{\textbf{[B]}}}
\newcommand{\Std}{\textcolor{standard}{\textbf{[S]}}}
\newcommand{\E}{\textcolor{erweiterung}{\textbf{[E]}}}

% ============================================================================
% DOKUMENT
% ============================================================================
\begin{document}

% ============================================================================
% DECKBLATT
% ============================================================================
\begin{titlepage}
    \centering
    \vspace*{2cm}
    {\large\textcolor{gray}{\schulart}}\\[2cm]
    {\Huge\textcolor{fach}{\textbf{Langentwurf}}}\\[0.5cm]
    {\large LE-\stundennummer{} | Tier 1 (vollständig)}\\[2cm]
    {\LARGE\textbf{\stundenthema}}\\[0.5cm]
    {\large Sequenz: \reihenthema{} -- \textbf{ABSCHLUSS}}\\[2cm]
    \begin{tabular}{rl}
        \textbf{Fach:} & \fach{} (A1) \\[0.3cm]
        \textbf{Klasse:} & \jahrgangsstufe \\[0.3cm]
        \textbf{Lehrwerk:} & \lehrwerk, \lehrwerkseite \\[0.3cm]
        \textbf{Doppelstunde:} & \stundennummer{} (Std. 7-8 von 8) \\[0.3cm]
    \end{tabular}
    \vfill
\end{titlepage}

\tableofcontents
\newpage

% ============================================================================
% 1. BEDINGUNGSANALYSE
% ============================================================================
\section{Bedingungsanalyse}

\subsection{Lerngruppe}
\begin{tabular}{ll}
    Kursgröße: & \lerngruppengroesse{} SuS \\
    Jahrgangsstufe: & \jahrgangsstufe \\
    Sprachniveau: & A1 (GER) -- Anfänger \\
    Fremdsprache: & 2. Fremdsprache (nach Englisch) \\
\end{tabular}

\subsection{Vorwissen aus Sequenz 2 (2a/2b/2c)}
\begin{itemize}
    \item \textbf{2a:} Familienmitglieder (madre, padre, hermano/a, abuelo/a, tío/a, primo/a)
    \item \textbf{2b:} Das Verb \textit{tener} (haben) -- Konjugation und Anwendung
    \item \textbf{2c:} Possessivbegleiter (mi, tu, su, nuestro/a) + Zahlen 1-20
\end{itemize}

\subsection{Differenzierung (intern)}
\begin{tabular}{|l|p{3.5cm}|p{3.5cm}|p{3.5cm}|}
\hline
& \B{} \textbf{Basis} & \Std{} \textbf{Standard} & \E{} \textbf{Erweiterung} \\
\hline
\textbf{Ziel} & BR: Rezeption, Grundwortschatz & MR: Produktion mit Hilfe & GY: Freie Anwendung \\
\hline
\textbf{Familia} & Zuordnung mit Bild & Sätze mit Wortschatz & Familienbeschreibung \\
\hline
\textbf{tener} & Formen einsetzen & Sätze bilden & Freie Anwendung \\
\hline
\textbf{Test} & Basisaufgaben (60\%) & Alle Pflichtaufgaben & + kreative Zusätze \\
\hline
\end{tabular}

% ============================================================================
% 2. SACHANALYSE
% ============================================================================
\section{Sachanalyse}

\subsection{Konsolidierung: Familienwortschatz}

\textbf{Kernwortschatz Familie:}
\begin{center}
\begin{tabular}{|l|l|l|}
\hline
\textbf{Spanisch} & \textbf{Deutsch} & \textbf{Plural} \\
\hline
la madre & die Mutter & las madres \\
el padre & der Vater & los padres \\
los padres & die Eltern & -- \\
la hermana & die Schwester & las hermanas \\
el hermano & der Bruder & los hermanos \\
los hermanos & die Geschwister & -- \\
la abuela & die Großmutter & las abuelas \\
el abuelo & der Großvater & los abuelos \\
los abuelos & die Großeltern & -- \\
la tía & die Tante & las tías \\
el tío & der Onkel & los tíos \\
la prima & die Cousine & las primas \\
el primo & der Cousin & los primos \\
\hline
\end{tabular}
\end{center}

\subsection{Konsolidierung: Das Verb tener}

\textbf{Konjugation von \textit{tener} (haben):}
\begin{center}
\begin{tabular}{|l|l|l|}
\hline
\textbf{Person} & \textbf{Form} & \textbf{Deutsch} \\
\hline
yo & tengo & ich habe \\
tú & tienes & du hast \\
él/ella/usted & tiene & er/sie hat; Sie haben \\
nosotros/as & tenemos & wir haben \\
vosotros/as & tenéis & ihr habt \\
ellos/ellas/ustedes & tienen & sie haben; Sie haben \\
\hline
\end{tabular}
\end{center}

\textbf{Besonderheit:} Stammwechsel e $\rightarrow$ ie (außer 1. Pers. Sg. und 1./2. Pers. Pl.)

\subsection{Konsolidierung: Possessivbegleiter}

\begin{center}
\begin{tabular}{|l|l|l|l|l|}
\hline
& \textbf{mask. Sg.} & \textbf{fem. Sg.} & \textbf{mask. Pl.} & \textbf{fem. Pl.} \\
\hline
mein/e & mi & mi & mis & mis \\
dein/e & tu & tu & tus & tus \\
sein/ihr/Ihr & su & su & sus & sus \\
unser/e & nuestro & nuestra & nuestros & nuestras \\
euer/eure & vuestro & vuestra & vuestros & vuestras \\
ihr/Ihr & su & su & sus & sus \\
\hline
\end{tabular}
\end{center}

\textbf{Wichtig:} Possessivbegleiter richten sich nach dem Besitzobjekt, nicht nach dem Besitzer!

\subsection{Konsolidierung: Zahlen 1-20}

\begin{center}
\begin{tabular}{|c|l|c|l|c|l|c|l|}
\hline
1 & uno & 6 & seis & 11 & once & 16 & dieciséis \\
2 & dos & 7 & siete & 12 & doce & 17 & diecisiete \\
3 & tres & 8 & ocho & 13 & trece & 18 & dieciocho \\
4 & cuatro & 9 & nueve & 14 & catorce & 19 & diecinueve \\
5 & cinco & 10 & diez & 15 & quince & 20 & veinte \\
\hline
\end{tabular}
\end{center}

% ============================================================================
% 3. DIDAKTISCHE ANALYSE
% ============================================================================
\section{Didaktische Analyse}

\subsection{Stellung in der Sequenz}
Dies ist die \textbf{4. und letzte Doppelstunde} (Std. 7-8 von 8) der Sequenz ``\reihenthema''.
Die Stunde schließt die Familiensequenz ab mit:
\begin{enumerate}
    \item Umfassende Wiederholung aller Inhalte (2a-2c)
    \item Übungsphase mit Arbeitsblatt
    \item Stundenleistung (summative Bewertung)
\end{enumerate}

\subsection{Kompetenzziele (Rahmenplan MV)}
\begin{itemize}
    \item \textbf{Sprechen:} Die eigene Familie beschreiben
    \item \textbf{Schreiben:} Kurze Texte über die Familie verfassen
    \item \textbf{Verfügen über sprachliche Mittel:} tener konjugieren, Possessivbegleiter anwenden
    \item \textbf{Interkulturell:} Familienstrukturen in Spanien/Lateinamerika
\end{itemize}

\subsection{Legitimation}
\textbf{Rahmenplan Spanisch MV:}
\begin{itemize}
    \item Zusammenhängendes Sprechen: über die eigene Familie sprechen
    \item Verfügen über sprachliche Mittel: Possessivbegleiter, Zahlen bis 20
    \item Schreiben: kurze persönliche Texte verfassen
\end{itemize}

% ============================================================================
% 4. STUNDENZIELE
% ============================================================================
\section{Stundenziele}

\subsection{Grobziel}
Die SuS \textbf{konsolidieren} alle Inhalte der Sequenz 2 (Familienwortschatz, tener, Possessivbegleiter, Zahlen) und weisen ihre Kompetenzen in einer \textbf{Stundenleistung} nach.

\subsection{Feinziele (operationalisiert)}
\begin{tabular}{|c|p{8cm}|c|c|}
\hline
\textbf{Nr.} & \textbf{Die SuS können...} & \textbf{AFB} & \textbf{Diff.} \\
\hline
FZ1 & die Familienmitglieder auf Spanisch benennen und zuordnen & I & alle \\
\hline
FZ2 & das Verb \textit{tener} in allen Personen korrekt konjugieren & I/II & alle \\
\hline
FZ3 & Possessivbegleiter passend zu Genus und Numerus einsetzen & II & alle \\
\hline
FZ4 & einen kurzen zusammenhängenden Text über ihre Familie schreiben & II/III & \Std{}/\E{} \\
\hline
FZ5 & die Zahlen 1-20 in Sätzen korrekt anwenden (z.B. Alter) & I/II & alle \\
\hline
\end{tabular}

% ============================================================================
% 5. METHODISCHE ANALYSE
% ============================================================================
\section{Methodische Analyse}

\subsection{Medien}
\begin{itemize}
    \item Tafel/Whiteboard
    \item Lehrwerk: \lehrwerk, \lehrwerkseite
    \item Arbeitsblatt: AB-02d.pdf (Wiederholungsübungen)
    \item Stundenleistung: SL-02d.pdf (Test)
    \item Musterlösung: SL-02d-ML.pdf
    \item Optional: LP-02d.html (digitale Wiederholung)
\end{itemize}

\subsection{Sozialformen}
\begin{itemize}
    \item Plenum: Wiederholung, Vorbereitung Test
    \item Partnerarbeit: Dialogübungen Familie
    \item Einzelarbeit: AB-Übungen, Stundenleistung
\end{itemize}

\subsection{Besonderheit: Stundenleistung}
\begin{itemize}
    \item Dauer: 25 Minuten
    \item Umfang: 20 Punkte (4 Aufgaben)
    \item Bewertung: Note nach Prozentskala
    \item Differenzierung: Basis 60\% = ausreichend
\end{itemize}

% ============================================================================
% 6. VERLAUFSPLANUNG
% ============================================================================
\section{Verlaufsplanung}

\textbf{1. Stunde (45 Min.) -- Schwerpunkt: Wiederholung + Übung}

\begin{longtable}{|p{1cm}|p{1.5cm}|p{4cm}|p{3cm}|p{2cm}|p{2cm}|}
\hline
\textbf{Zeit} & \textbf{Phase} & \textbf{Lehrerhandeln} & \textbf{Schülerhandeln} & \textbf{SF} & \textbf{Medien} \\
\hline
\endhead

0--5 & Einstieg & Begrüßung, Ankündigung: Wiederholung + Test & Zuhören & Plenum & -- \\
\hline

5--12 & Wieder\-holung 1 & Familienwortschatz abfragen: ``¿Cómo se dice Mutter?'' & Antworten: ``la madre'' & Plenum & Tafel \\
\hline

12--20 & Wieder\-holung 2 & tener-Konjugation: Formen an die Tafel & Mitsprechen, Notieren & Plenum & Tafel \\
\hline

20--30 & Übung AB & AB-02d ausgeben, Aufgaben 1-2 (Zuordnung, tener) & Bearbeiten & EA & AB \\
\hline

30--40 & Übung AB & AB-02d Aufgaben 3-4 (Possessiv, Text) & Bearbeiten & EA & AB \\
\hline

40--45 & Sicherung & Lösungen besprechen, Fragen klären & Vergleichen, Fragen & Plenum & Tafel \\
\hline
\end{longtable}

\vspace{1em}

\textbf{2. Stunde (45 Min.) -- Schwerpunkt: Stundenleistung}

\begin{longtable}{|p{1cm}|p{1.5cm}|p{4cm}|p{3cm}|p{2cm}|p{2cm}|}
\hline
\textbf{Zeit} & \textbf{Phase} & \textbf{Lehrerhandeln} & \textbf{Schülerhandeln} & \textbf{SF} & \textbf{Medien} \\
\hline
\endhead

0--5 & Aktivierung & Letzte Wiederholung: ``Tengo una hermana'' etc. & Beispielsätze bilden & Plenum & -- \\
\hline

5--10 & Vorbereitung & SL-02d erklären, Fragen klären, Ruhe einkehren & Zuhören, Fragen & Plenum & SL \\
\hline

10--35 & Test & Aufsicht, individuelle Hilfe bei Verständnisfragen & Stundenleistung bearbeiten & EA & SL \\
\hline

35--40 & Abgabe & Einsammeln, kurze Entspannung & Abgeben & -- & -- \\
\hline

40--45 & Abschluss & Ausblick Sequenz 3 (Mi casa), Feedback & Rückmeldung & Plenum & -- \\
\hline
\end{longtable}

\textbf{Legende:} EA = Einzelarbeit, PA = Partnerarbeit, SF = Sozialform, SL = Stundenleistung

% ============================================================================
% 7. ERWARTETE SCHWIERIGKEITEN
% ============================================================================
\section{Erwartete Schwierigkeiten}

\begin{tabular}{|p{5cm}|p{8cm}|}
\hline
\textbf{Schwierigkeit} & \textbf{Reaktion} \\
\hline
Verwechslung tener/ser & Bedeutungsunterschied klären: tener = haben, ser = sein \\
\hline
Genus bei Possessivbegleitern & Regel wiederholen: richtet sich nach dem Objekt! \\
\hline
Prüfungsangst bei Stundenleistung & Ruhige Atmosphäre, klare Instruktionen, Zeitzugabe für \B{} \\
\hline
Zeitdruck bei SL & Klare Zeitansagen, einfachere Aufgaben zuerst \\
\hline
Schreibfehler bei Familienwörtern & Akzeptieren, wenn phonetisch korrekt (z.B. ``ermano'') \\
\hline
\end{tabular}

% ============================================================================
% 8. HAUSAUFGABE
% ============================================================================
\section{Hausaufgabe}

\begin{itemize}
    \item \textbf{Vor der Stunde:} LP-02d.html zur Wiederholung (optional)
    \item \textbf{Nach der Stunde:} Vokabeln Sequenz 2 wiederholen für Vokabeltest
\end{itemize}

% ============================================================================
% 9. LÖSUNGEN AB-02d
% ============================================================================
\section{Lösungen AB-02d}

\textbf{Aufgabe 1: Familienmitglieder zuordnen} (4P)
\begin{itemize}
    \item la madre $\rightarrow$ die Mutter
    \item el padre $\rightarrow$ der Vater
    \item la hermana $\rightarrow$ die Schwester
    \item el abuelo $\rightarrow$ der Großvater
\end{itemize}

\textbf{Aufgabe 2: tener konjugieren} (4P)
\begin{itemize}
    \item yo tengo
    \item tú tienes
    \item ella tiene
    \item nosotros tenemos
\end{itemize}

\textbf{Aufgabe 3: Possessivbegleiter einsetzen} (6P)
\begin{itemize}
    \item mi madre (meine Mutter)
    \item tu hermano (dein Bruder)
    \item su padre (sein/ihr Vater)
    \item nuestra abuela (unsere Großmutter)
    \item mis hermanos (meine Geschwister)
    \item sus primos (seine/ihre Cousins)
\end{itemize}

\textbf{Aufgabe 4: Text über Familie schreiben} (6P)

Beispiellösung:
\begin{quote}
Mi familia es grande. Tengo una madre y un padre. Mi madre se llama María y tiene cuarenta años. También tengo dos hermanos. Mi hermano se llama Pablo y mi hermana se llama Ana. Mis abuelos viven en Madrid.
\end{quote}

Bewertung:
\begin{itemize}
    \item 3+ Familienmitglieder genannt: 2P
    \item tener korrekt verwendet: 2P
    \item Possessivbegleiter korrekt: 1P
    \item Zusammenhängender Text: 1P
\end{itemize}

% ============================================================================
% 10. STUNDENLEISTUNG - ÜBERSICHT
% ============================================================================
\section{Stundenleistung SL-02d}

\textbf{Aufbau:} 4 Aufgaben, 20 Punkte, 25 Minuten

\begin{tabular}{|c|l|c|c|c|}
\hline
\textbf{Aufg.} & \textbf{Inhalt} & \textbf{AFB} & \textbf{Punkte} & \textbf{Diff.} \\
\hline
1 & Zuordnung: Familienmitglieder & I & 4 & alle \\
\hline
2 & Lückentext: tener einsetzen & I/II & 4 & alle \\
\hline
3 & Possessivbegleiter + Zahlen & II & 6 & alle \\
\hline
4 & Dialog vervollständigen & II/III & 6 & alle \\
\hline
\multicolumn{3}{|r|}{\textbf{Gesamt:}} & \textbf{20} & \\
\hline
\end{tabular}

\vspace{1em}

\textbf{Bewertung (Klasse 7):}
\begin{tabular}{|c|c|c|c|c|c|c|}
\hline
Note & 1 & 2 & 3 & 4 & 5 & 6 \\
\hline
ab \% & 96 & 80 & 60 & 40 & 20 & <20 \\
\hline
ab Punkte & 19 & 16 & 12 & 8 & 4 & <4 \\
\hline
\end{tabular}

% ============================================================================
% 11. DIDAKTIK-SCORE
% ============================================================================
\newpage
\section{Didaktik-Score}

\begin{scorebox}
\textbf{Bewertung der didaktischen Qualität nach 10 Dimensionen}\\
\textbf{Ziel-Score: $\geq$ 9.0 (Premium-Qualität)}

\vspace{1em}
\begin{tabular}{|c|l|c|c|p{4.5cm}|}
\hline
\textbf{\#} & \textbf{Dimension} & \textbf{Gew.} & \textbf{Score} & \textbf{Notiz} \\
\hline
1 & Differenzierung & 15\% & 9/10 & [B]/[S]/[E] durchgängig, Zeitzugabe \\
\hline
2 & Taxonomie (AFB) & 10\% & 9/10 & AFB I: 35\%, AFB II: 45\%, AFB III: 20\% \\
\hline
3 & Lernziel-Alignment & 10\% & 10/10 & Rahmenplan: Familie beschreiben \\
\hline
4 & Aktivierung & 10\% & 9/10 & Wiederholung interaktiv, PA-Dialoge \\
\hline
5 & Scaffolding & 10\% & 9/10 & Wiederholung $\rightarrow$ Übung $\rightarrow$ Test \\
\hline
6 & Feedback-Qualität & 10\% & 9/10 & Bewertungsraster transparent, Lösungen \\
\hline
7 & Sprachliche Klarheit & 10\% & 10/10 & Altersgerecht Kl. 7 \\
\hline
8 & Motivation \& Relevanz & 10\% & 9/10 & Familie = Alltagsrelevanz \\
\hline
9 & Inklusion (UDL) & 10\% & 9/10 & Zeitzugabe, mündlich/schriftlich \\
\hline
10 & Praktikabilität & 5\% & 10/10 & 45+45 Min realistisch, klare Struktur \\
\hline
\multicolumn{3}{|r|}{\textbf{GESAMT:}} & \textbf{9.2/10} & \textcolor{successgreen}{\textbf{FREIGABE}} \\
\hline
\end{tabular}
\end{scorebox}

\end{document}
